\chapter{Bandha and Mudra}
\label{ch:bandha-mudra}

\emph{What are the locks and seals of yoga, and which ones must the candidate know?}
The syllabus lists three bandhas and a set of hasta (hand) mudras. The two are separate topics in the syllabus, though the Hatha Pradipika treats them together in one chapter. Bandhas are learned with a teacher and asked about mostly in theory; the hand mudras are used in every pranayama and meditation the candidate demonstrates.

\begin{unverified}
This chapter is written from notes blocks whose fact check is still pending. The bandhas follow the Hatha Pradipika (an English translation on Wikisource) and the Board syllabus; jnana and nasagra mudra follow the Common Yoga Protocol; the other hand mudras come from the workshop slides and still need a named source. A certified teacher must review it before print.
\end{unverified}

\section{Bandha}

\subsection*{Meaning}
A bandha is a lock: a gentle muscular hold that is said to direct prana and steady the mind. Bandhas are done with pranayama, after the asanas. The tradition says they guide prana and bring stability and focus; take that as the tradition's own account, not as a medical claim. Link each lock to its place in the body: the throat, the belly, and the base of the trunk.

\subsection*{Jalandhara bandha (throat lock)}
Sit tall. Breathe in and hold the breath; lower the chin to the chest and hold. Then raise the head and breathe out. The Hatha Pradipika describes it as contracting the throat and pressing the chin firmly on the chest. The guidebook shows the same chin lock in the shoulder stand.

\subsection*{Uddiyana bandha (abdominal lock)}
Done standing or seated, and only on an empty stomach. Breathe out fully; draw the belly in and up under the ribs, and hold without breath. Release slowly, then breathe in. The Hatha Pradipika describes it as drawing the belly, above and below the navel, back towards the spine.

\subsection*{Mula bandha (root lock)}
Contract the muscles of the perineum, the pelvic floor, and lift gently. Hold with or without the breath and release slowly. In the Hatha Pradipika the heel presses the perineum and the lower breath (apana) is drawn upward by contracting the anus.

\subsection*{Maha bandha}
The three locks done together are the great lock. The Hatha Pradipika groups it with the great mudra and the great piercing, and says that all the bandhas are to be learned from a teacher and practised little by little. It is taught only to practised students. The text's claims for them, such as rejuvenation and victory over death, are tradition and are not health claims.

Schools differ on whether a lock is held after the in-breath or after the out-breath. Teach the Board's form.

\begin{caution}
Do not practise the bandhas in pregnancy, high blood pressure, heart disease, hernia, or a stomach or intestinal ulcer; avoid uddiyana and mula bandha during menstruation. Learn them only with a qualified teacher, and never force.
\end{caution}

\ycbfig{ycb1-fig-bandhas}

\section{Mudra}

\subsection*{Meaning and types}
A mudra is a seal or gesture, used with asana, pranayama and meditation. A traditional gloss reads the word as ``that which gives delight'' (\emph{mudam ānandam dadāti iti mudrā}). The traditional idea is that the hand and finger positions keep the energy of prana turned back into the body. The workshop slides group mudras into five families:
\begin{itemize}
\item \textbf{Hasta}: hand gestures, the ones in the syllabus.
\item \textbf{Kaya}: whole-body gestures.
\item \textbf{Mukha}: gestures of the face and eyes, such as khechari (a tongue gesture), bhrumadhya and shambhavi (gazing at the eyebrow centre).
\item \textbf{Bandha}: the internal locks above.
\item \textbf{Chakra}: gestures named after the centres, such as muladhara, anahata and ajna.
\end{itemize}
This grouping is a school classification; the Board lists bandha and mudra as separate topics.

Teach every mudra as a position to keep relaxed and effortless. The fingers must not become tense.

\subsection*{Jnana mudra}
The tip of the thumb touches the tip of the index finger to form a circle; the other three fingers stay straight, relaxed and together. The Common Yoga Protocol holds it with the palms up on the thighs. It is the hand of meditation and the resting hand in pranayama.

\subsection*{Chin mudra}
The Board syllabus lists chin and jnana mudra separately, but the sources read do not state the difference. Many popular texts give chin mudra with the palms up and jnana with the palms down; others use the two names for the same gesture. Confirm the Board's form with a certified teacher before teaching a difference.

\subsection*{Nasagra mudra}
The hand used to close the nostrils in nadi shodhana: the index and middle fingers folded, the thumb by the right nostril, the ring and little fingers by the left. It is also called nasika or Vishnu mudra. Some slides describe it with the thumb and ring finger only; that does not match the protocol.

\subsection*{The other hand mudras of the syllabus}
\textit{Source: workshop slides, to be checked.}
\begin{itemize}
\item \textbf{Chinmaya mudra}: the index tip touches the thumb tip and the other fingers curl into the palm like a loose fist, palms up.
\item \textbf{Brahma mudra}: the thumb is folded inside towards the little finger and the other fingers close over it into a fist; the two fists are joined and held at the navel.
\item \textbf{Adi mudra}: the same fist as brahma mudra, with the hands kept apart on the thighs.
\item \textbf{Dhyana mudra}: both palms up, one resting on the other, in the lap below the navel.
\end{itemize}

\ycbfig{ycb1-fig-hastamudra}
\ycbfig{ycb1-fig-kayamudra}

\begin{recall}
\begin{itemize}
\item Bandhas: jalandhara (throat), uddiyana (belly), mula (base); all three together is maha bandha.
\item Uddiyana only on an empty stomach, after a full out-breath.
\item Bandhas are learned with a teacher; many cautions, including pregnancy, high blood pressure and heart disease.
\item Jnana mudra: thumb and index tips joined, other fingers straight; palms up in the protocol.
\item Nasagra mudra: index and middle folded; used in nadi shodhana.
\end{itemize}
\end{recall}

\subsection*{Worked examples}
\begin{itemize}
\item \textbf{A student who has just eaten asks to try uddiyana.} Decline: uddiyana is done only on an empty stomach. Offer gentle breathing instead.
\item \textbf{A student's fingers grip tightly in jnana mudra.} Ask for a lighter touch, with the three fingers straight but soft; the mudra should be effortless.
\item \textbf{Model answer: name each bandha with its place.} Jalandhara at the throat, uddiyana at the abdomen, mula at the perineum.
\end{itemize}

\subsection*{Check yourself}
\begin{enumerate}
\item Name the three bandhas in the syllabus.
\item Which bandha is the throat lock? The root lock? The abdominal lock?
\item On what condition is uddiyana done?
\item What is maha bandha?
\item Give three cautions for the bandhas.
\item Describe jnana mudra.
\item Which hand gesture is used with alternate nostril breathing?
\item What is adi mudra?
\item Name the five families of mudra used on the workshop slides.
\end{enumerate}

\subsection*{Answers}
\begin{enumerate}
\item Jalandhara, uddiyana and mula bandha.
\item Jalandhara; mula; uddiyana.
\item On an empty stomach, after a full breath out.
\item The three locks done together.
\item Any three of: pregnancy, high blood pressure, heart disease, hernia, ulcer, menstruation (for uddiyana and mula).
\item Thumb tip and index tip joined in a circle, the other three fingers straight, relaxed and together; palms up on the thighs in the protocol.
\item Nasagra (nasika or Vishnu) mudra.
\item Fists with the thumb inside, kept apart on the thighs.
\item Hasta, kaya, mukha, bandha and chakra mudras.
\end{enumerate}

\paragraph{Sources.} Hatha Pradipika of Svatmarama, third chapter (English translation on Wikisource, 1915); the Board syllabus for Level 1; the Common Yoga Protocol, Ministry of AYUSH (jnana and nasagra mudra); the official guidebook for Level I (2016, older scheme); the workshop slides for the remaining hand mudras (to be replaced by a named source).
